% Fichas de pesquisa: MER, CLI, DRE, ENG
% Conferido na web em 23/09/2026. "(a confirmar)" marca o que não foi verificado em fonte.

\section{Mercado, clima, resposta da demanda e engine}

% ---------------------------------------------------------------------------
\subsection{MER \textperiodcentered{} Mercado e preço}

\begin{nota}
\textbf{Trilha condicional.} As três features deste módulo, com CAR-08 e
GER-08, estão fora do roadmap faseado: entram apenas se a questão de
cliente-alvo for resolvida a favor de gerador e comercializadora. A pesquisa
fica registrada para que a decisão, quando vier, não recomece do zero.
\end{nota}

\begin{nota}
\textbf{Correção sobre a CCEE.} O portal Dados Abertos da CCEE tem um bloqueio
de segurança (firewall de aplicação) que barra clientes fora da política de
acesso e orienta a abrir chamado com o código do erro e o IP [1,2]. A própria
CCEE divulga tutoriais de automação por API, e o download de um endpoint
CKAN (\texttt{/datastore/dump/<id>}) funcionou sem bloqueio. Coleta
automatizada é viável com boas práticas: poucas requisições, user-agent
identificável, cache local e download dos CSVs anuais em vez de raspagem.
\end{nota}

\begin{ficha}{MER-01 \textperiodcentered{} PLD previsto por submercado (quantis)}{V3}
\campo{Dados} \emph{PLD\_HORARIO} da CCEE: um CSV por ano, horário de 2021 a
2026 e semanal de 2001 a 2020, licença CC-BY 4.0, atualização diária [3].
Limites regulatórios de 2026: piso R\$~57,31/MWh, teto estrutural
R\$~785,27/MWh e teto horário R\$~1.611,04/MWh (Despacho ANEEL 3.850/2025)
[5]. Decks NEWAVE/DECOMP/DESSEM no Acervo CCEE, inclusive decks sombra [6];
decks do DESSEM também no ONS [7]. ONS Dados Abertos: carga semi-horária,
geração por usina e CMO semi-horário [8].
\campo{Abordagens} Linha de base: ingênuo sazonal e LEAR (autorregressivo
com LASSO), referência da literatura, implementado no \emph{epftoolbox} [10].
Recomendada: gradient boosting com perda quantílica sobre CMO, carga e
renováveis previstas, armazenamento dos reservatórios e deck vigente, com
calibração conformal (a confirmar). Estado da arte: emulação do DESSEM por
modelo substituto e redes que preveem distribuições paramétricas [10].
\campo{Métricas} Pinball e CRPS, cobertura e largura dos intervalos, MAE
relativo ao ingênuo, teste de Diebold--Mariano.
\campo{Armadilhas} Quebra estrutural em 01/01/2021 (PLD passou a horário,
calculado pelo DESSEM) [4]; nova função de custo futuro em 12/10/2025 [6];
distribuição censurada no piso e no teto; séries do ONS revisadas depois de
publicadas, o que exige guardar o que era conhecido em cada data [8]. O
Projeto Meta II propõe preço por oferta e dois PLDs (ex-ante e ex-post), com
implantação prevista para junho de 2028 [9]: deriva estrutural garantida.
\end{ficha}

\begin{ficha}{MER-02 \textperiodcentered{} Cenários conjuntos de geração, carga e preço}{V3}
\campo{Dados} Saídas probabilísticas de GER-01, CAR-08 e MER-01; histórico
conjunto do ONS e da CCEE [3,8].
\campo{Abordagens} Linha de base: cenários independentes ou bootstrap de
blocos históricos. Recomendada: cópula gaussiana sobre as marginais
preditivas (Pinson et al.\ 2009) [11]; alternativa com a ordem dos membros do
ensemble meteorológico (ECC e dual-ECC) [12]. Estado da arte: cópulas vine e
geradores profundos (a confirmar).
\campo{Métricas} Energy score e variogram score, histogramas de posto
multivariados, CRPS por marginal (biblioteca \emph{scoringrules}) [13].
\campo{Armadilhas} A dependência muda com o regime hidrológico, e uma cópula
estacionária subestima eventos conjuntos extremos (seca com pico de carga).
\end{ficha}

\begin{ficha}{MER-03 \textperiodcentered{} Otimização de oferta e despacho de baterias}{V4}
\campo{Dados} Cenários de MER-02, PLD horário e parâmetros técnicos da
bateria. Contexto: Portaria MME 136/2026 define o primeiro leilão de reserva
de capacidade em baterias (dezembro de 2026, contratos de 15 anos a partir
de agosto de 2028), com receita de disponibilidade e não de arbitragem livre
[14].
\campo{Abordagens} Linha de base: otimização determinística sobre a previsão
pontual. Recomendada: programação estocástica em dois estágios ou CVaR,
em janela móvel. Estado da arte: aprendizado focado na decisão; minimizar o
CRPS não garante a melhor oferta [15].
\campo{Métricas} Receita contra o ótimo com informação perfeita (regret),
CVaR da receita, custo de degradação.
\campo{Armadilhas} Hoje o PLD vem de modelo de custo, não de ofertas; ``oferta''
só faz sentido com o modelo híbrido (2028 em diante) [9]. Degradação
simplificada superestima receita. Baixo encaixe com o mapa.
\end{ficha}

% ---------------------------------------------------------------------------
\subsection{CLI \textperiodcentered{} Contexto climático}

\begin{ficha}{CLI-01 \textperiodcentered{} Tempo previsto}{V2}
\campo{Dados} ECMWF Open Data: IFS e AIFS (determinístico e ensemble), grade
de 0,25°, quatro rodadas por dia, licença CC-BY 4.0 com uso comercial
permitido e atribuição obrigatória; o portal guarda só as rodadas dos
últimos 2 a 3 dias, com espelhos em nuvem [16,17]. NOAA GFS 0,25° no AWS,
aberto, com histórico longo no NCAR [18]. CPTEC/INPE: BAM, WRF, Eta e BRAMS,
com citação exigida [19]. INMET: estações automáticas horárias (portal com
90 dias) e BDMEP para o histórico, com defasagem de 1 a 3 meses [20].
\campo{Abordagens} Linha de base: previsão bruta interpolada para as
entidades. Recomendada: pós-processamento estatístico por entidade (MOS,
EMOS, regressão quantílica) com downscaling. Estado da arte: modelos globais
de ML (AIFS, GraphCast-GFS) [16,18]. Ferramentas: \emph{herbie},
\emph{xarray}/\emph{cfgrib}, \emph{xskillscore} [17].
\campo{Métricas} RMSE e MAE por variável e horizonte, CRPS do ensemble,
diagrama de confiabilidade, skill contra climatologia e persistência.
\campo{Armadilhas} \textbf{Arquivar cada rodada desde o primeiro dia}: sem
isso, não há histórico de previsões para treinar sem vazamento. Mudanças de
versão dos modelos quebram o pós-processamento. A atribuição do ECMWF precisa
aparecer no mapa.
\end{ficha}

\begin{ficha}{CLI-02 \textperiodcentered{} Satélite e radar observados}{V2}
\campo{Dados} GOES-East é o \textbf{GOES-19} desde abril de 2025 (75,2°W),
com ABI e GLM no AWS Open Data em NetCDF4; houve interrupções (ABI restaurado
em julho de 2026) [21]. Radar REDEMET/DECEA por API com chave, mas em imagem
PNG, útil para exibir e não para modelar [22]. CEMADEN: pluviômetros e radares
de dupla polarização, com dados volumétricos sob solicitação (a confirmar)
[23].
\campo{Abordagens} Linha de base: exibir mosaicos. Recomendada: nowcasting de
precipitação de 0 a 6 h com \emph{pysteps} (fluxo óptico e ensembles
estocásticos) [24]. Estado da arte: nowcasting profundo (ConvGRU, difusão).
\campo{Métricas} CSI, POD e FAR por limiar; FSS; CRPS [24].
\campo{Armadilhas} Cobertura de radar desigual no território; lacunas e
latência do GOES.
\end{ficha}

\begin{ficha}{CLI-03 \textperiodcentered{} Raios e eventos severos}{V3}
\campo{Dados} BrasilDAT (ELAT/INPE): detecta descargas nuvem-solo e
intranuvem; acesso por contato com o ELAT, \textbf{não é aberta} [25].
Alternativa aberta: GLM do GOES-19, detecção óptica de raios totais com
eficiência uniforme sobre o Brasil [21,26].
\campo{Abordagens} Densidade de descargas por célula e janela; cruzamento
espaço-temporal de raios com desligamentos de alimentadores (alimenta
CON-03).
\campo{Métricas} POD e FAR dos alertas contra desligamentos; antecedência.
\campo{Armadilhas} O GLM registra flashes ópticos, não pontos de impacto, com
localização da ordem de quilômetros (a confirmar): não serve para associar
descarga a vão de linha.
\end{ficha}

% ---------------------------------------------------------------------------
\subsection{DRE \textperiodcentered{} Resposta da demanda e consumidor}

\begin{ficha}{DRE-01 \textperiodcentered{} Alvos de resposta da demanda por alimentador}{V4}
\campo{Dados} Regulação: REN ANEEL 1.040/2022 (programa estrutural) e sandbox
do ONS para o produto disponibilidade, com primeiro mecanismo competitivo em
outubro de 2024 (até 4 acionamentos mensais de 4 h, entre 18 h e 22 h)
[27,28]. Resultados na CCEE [28]. O programa do ONS mira grandes consumidores
do SIN; resposta da demanda por alimentador é caso da distribuidora e depende
de medição própria (a confirmar).
\campo{Abordagens} Linha de base de medição: regras ``X de Y dias'' (a
confirmar). Recomendada: baseline contrafactual por ML, avaliada como
contrafactual porque define pagamento; estado da arte: controle sintético
generalizado e baselines probabilísticos [29]. Alvos: alimentadores com
carregamento previsto acima do limite (CAR-01, RED-01), ordenados por alívio
esperado por custo.
\campo{Métricas} Viés da baseline em dias sem evento; precisão no topo dos
alimentadores escolhidos; MW reduzidos contra contratados.
\campo{Armadilhas} Baseline enviesada paga indevidamente ou incentiva inflar
a carga antes do evento; efeito rebote; poucos eventos para validar.
\end{ficha}

\begin{ficha}{DRE-02 \textperiodcentered{} Desagregação de consumo (NILM)}{fora}
\campo{Dados} Medição agregada de alta frequência por unidade; bases públicas
REDD e UK-DALE; nenhuma base brasileira encontrada [30].
\campo{Abordagens} FHMM (NILMTK) como linha de base; seq2point (Zhang et
al.\ 2018) como recomendada [30].
\campo{Por que fica fora} Exige dado individual de alta frequência que não
existe em escala, o parque de aparelhos brasileiro (chuveiro elétrico) difere
das bases de treino, e a saída é por unidade consumidora, não espacial.
\end{ficha}

\begin{ficha}{DRE-03 \textperiodcentered{} Recomendação tarifária e de eficiência}{fora}
\campo{Dados} Tarifa Branca desde 2018, adesão de cerca de 69 mil unidades em
75 milhões aptas [31,32]; proposta da área técnica da ANEEL (novembro de 2025)
de Tarifa Branca automática para consumo acima de 1 MWh/mês, levada a consulta
pública, desfecho não verificado [32]. Tarifas homologadas por distribuidora
nos dados abertos da ANEEL [33].
\campo{Por que fica fora} Sem curva horária medida a recomendação é
especulativa, e a decisão é da unidade consumidora, não de uma entidade do
mapa. Se a Tarifa Branca automática vingar, a versão agregada (quantas
unidades de um alimentador seriam afetadas) pode voltar como insight.
\end{ficha}

% ---------------------------------------------------------------------------
\subsection{ENG \textperiodcentered{} Engine de análise e modelos}

\begin{ficha}{ENG-01 \textperiodcentered{} Feed de insights consolidado}{MVP}
\campo{Dados} Todas as camadas com a hierarquia, uma matriz de vizinhança
(geográfica ou por topologia elétrica) e o histórico por ciclo.
\campo{Abordagens} \emph{Pontuação}: adotar o esquema de Top-K Insights
(SIGMOD 2017) e QuickInsights (SIGMOD 2019, base do recurso do Power BI):
tipo de insight, subespaço e pontuação por impacto × significância [34].
\emph{Anomalia espacial}: Moran local (LISA) e Getis--Ord $G_i^*$ em
PySAL/\emph{esda}, com pseudo-p por permutação [35], sempre com controle de
FDR para testes múltiplos e dependentes [36]. \emph{Mudança e tendência}:
pontos de mudança offline com \emph{ruptures} e online com BOCPD [37];
excedência de quantil da previsão para limiares. \emph{Consolidação}: o
insight do pai absorve os dos filhos quando explica a maior parte do efeito;
a reconciliação MinT garante números coerentes entre níveis [38].
\campo{Métricas} Precisão no topo do feed (insights marcados úteis ou
acionados), taxa de falsos positivos sob dados embaralhados, insights por
evento real.
\campo{Armadilhas} Explosão de testes (entidades × camadas × detectores);
MAUP (o resultado muda com a escala de agregação); vizinhança geográfica não
é conectividade elétrica; fadiga de alertas.
\end{ficha}

\begin{ficha}{ENG-02 \textperiodcentered{} Explicação por entidade}{MVP}
\campo{Abordagens} TreeSHAP para modelos de árvore [39]; TimeSHAP para
modelos de séries [40]; GeoShapley para modelos espaciais, que trata a
localização como um jogador conjunto [41].
\campo{Métricas} Fidelidade (remoção e inserção), estabilidade entre
execuções, consistência com a física (sinal esperado da temperatura sobre a
carga).
\campo{Armadilhas} Variáveis climáticas correlacionadas dão atribuições
diferentes conforme a variante de SHAP [39]; explicação não é causa e não
deve alimentar a simulação de cenário como se fosse efeito.
\end{ficha}

\begin{ficha}{ENG-03 \textperiodcentered{} Erro dos modelos por entidade}{MVP}
\campo{Dados} Pares previsto × realizado, versionados por data de emissão,
horizonte e versão do modelo; realizados revisados pelo ONS e pela CCEE
exigem guardar a versão conhecida em cada data [8].
\campo{Métricas} nMAE com normalização explícita; CRPS (regra estritamente
própria) [42]; pinball; calibração (PIT, cobertura, confiabilidade) e
nitidez; precisão no topo para rankings. Bibliotecas: \emph{scoringrules}
(ativa, univariada e multivariada) e \emph{xskillscore}; a
\emph{properscoring} está pouco mantida [42].
\campo{Armadilhas} nMAE explode em entidades de carga baixa; métricas
agregadas escondem a cauda; poucas amostras por entidade pedem intervalo por
bootstrap antes de emitir o insight de desempenho.
\end{ficha}

\begin{ficha}{ENG-04 \textperiodcentered{} Deriva de dados por região}{V2}
\campo{Abordagens} Testes univariados (KS, PSI, Wasserstein) como no
\emph{Evidently} [43]; MMD para deriva multivariada; estimativa de desempenho
sem rótulo com NannyML [45]; testes de duas amostras sobre representações
reduzidas funcionam melhor segundo Rabanser et al.\ (2019) [46]. Aplicar o
teste por região e procurar agrupamentos de deriva com LISA (proposta, a
confirmar como prática).
\campo{Armadilhas} \textbf{Licença}: o \emph{alibi-detect} mudou de Apache 2.0
para Business Source License 1.1 em janeiro de 2024, com restrição a uso
comercial em produção [44]. Com amostras grandes tudo ``deriva'': usar
tamanho de efeito e FDR; comparar com a mesma estação do ano.
\end{ficha}

\begin{ficha}{ENG-05 \textperiodcentered{} Simulação de cenário}{V2}
\campo{Abordagens} Perturbar entradas e reexecutar os modelos, alterando de
forma coerente variáveis correlacionadas e sinalizando extrapolação fora do
domínio de treino. Onde houver cenários conjuntos (MER-02, na trilha
condicional), propagar a incerteza sobre eles; sem eles, perturbar as
marginais e declarar que a dependência não está modelada.
Estado da arte: modelos causais estruturais (a confirmar: DoWhy, EconML).
\campo{Métricas} Backtest em choques históricos conhecidos; cobertura dos
intervalos.
\campo{Armadilhas} Modelos preditivos não são causais; o usuário pode ler o
cenário como previsão.
\end{ficha}

\begin{ficha}{ENG-06 \textperiodcentered{} Pergunta em texto livre}{V2}
\campo{Estado da técnica} Em text-to-SQL sobre bancos reais (BIRD), o melhor
sistema chega a cerca de 82\% de acurácia de execução contra 93\% de humanos
[47]. Em tarefas geoespaciais, agentes frequentemente ``alucinam'' relações
espaciais em vez de computá-las [48]. No DataGemma, mesmo com recuperação de
dados, parte relevante das estatísticas citadas saiu errada [49]: o modelo de
linguagem não pode ser fonte de número.
\campo{Arquitetura recomendada} (1) o modelo escolhe ferramentas de uma API
semântica restrita (métricas, entidades, filtros); SQL livre só como recurso
de exceção, em leitura e com limites; (2) a engine executa e devolve números
com identificadores; (3) o texto é redigido com marcadores preenchidos pela
engine; (4) um verificador rejeita a resposta se algum número não vier do
resultado; (5) a resposta cita os insights e consultas e admite ``não sei''.
\campo{Métricas} Acurácia de execução num conjunto próprio de perguntas do
setor; taxa de números sem fonte (meta: zero); taxa de recusa correta.
\campo{Armadilhas} Unidade errada (MW e MWh); nomes de subestação repetidos;
injeção de instruções vinda de texto nos dados; acesso a dados de unidades
consumidoras (controle por papel na camada de ferramentas).
\end{ficha}

\subsection*{Referências desta seção}
{\small
\begin{enumerate}[label={[\arabic*]},leftmargin=2.2em,itemsep=1pt]
\item CCEE Dados Abertos, mensagem de bloqueio: \url{https://dadosabertos.ccee.org.br/}
\item CCEE Dados Abertos, sobre e contatos: \url{https://dadosabertos.ccee.org.br/about}
\item CCEE, PLD\_HORARIO: \url{https://dadosabertos.ccee.org.br/dataset/pld_horario}
\item CCEE, conceitos de preços: \url{https://www.ccee.org.br/precos/conceitos-precos}
\item ANEEL, PLD 2026: \url{https://www.gov.br/aneel/pt-br/assuntos/noticias/2025/aneel-define-tarifas-de-energia-de-otimizacao-de-servicos-ancilares-e-pld-para-2026}
\item CCEE, Acervo e novas funções de custo futuro: \url{https://www.ccee.org.br/acervo-ccee}
\item ONS, decks do DESSEM: \url{https://www.ons.org.br/}
\item ONS Dados Abertos: \url{https://dados.ons.org.br/}
\item CCEE, Projeto Meta II: \url{https://www.ccee.org.br/en/web/guest/-/ccee-apresenta-primeiros-resultados-do-projeto-meta-ii-sobre-formacao-de-preco-no-brasil}
\item Lago et al.\ (2021), \emph{Applied Energy} 293:116983; \url{https://github.com/jeslago/epftoolbox}
\item Pinson et al.\ (2009), \emph{Wind Energy} 12:51--62, doi:10.1002/we.284
\item Ben Bouallègue et al.\ (2016), dual-ECC, \emph{Monthly Weather Review} 144(12)
\item scoringrules: \url{https://scoringrules.readthedocs.io/}
\item MME, diretrizes do leilão de baterias: \url{https://www.gov.br/mme/}
\item Decision-focused predict-then-bid: \url{https://arxiv.org/abs/2505.01551}; CRPS e decisão: \url{https://arxiv.org/abs/2308.15443}
\item ECMWF Open Data: \url{https://www.ecmwf.int/en/forecasts/datasets/open-data}
\item Herbie: \url{https://herbie.readthedocs.io/}
\item GFS no AWS: \url{https://registry.opendata.aws/noaa-gfs-pds/}; NCAR d084001
\item CPTEC/INPE: \url{https://previsaonumerica.cptec.inpe.br/}
\item INMET BDMEP: \url{https://bdmep.inmet.gov.br/}
\item GOES no AWS: \url{https://registry.opendata.aws/noaa-goes/}; NCEI ABI/GLM
\item API REDEMET, radar: \url{https://ajuda.decea.mil.br/base-de-conhecimento/api-redemet-produtos-radar/}
\item CEMADEN, mapa interativo: \url{https://mapainterativo.cemaden.gov.br/}
\item Pulkkinen et al.\ (2019), pysteps, \emph{GMD} 12:4185
\item ELAT/INPE, BrasilDAT: \url{https://www.gov.br/inpe/}
\item GLM, eficiência de detecção: \url{https://ntrs.nasa.gov/citations/20180000609}
\item REN ANEEL 1.040/2022: \url{https://www2.aneel.gov.br/cedoc/ren20221040.html}
\item CCEE e ONS, sandbox de resposta da demanda: \url{https://www.ccee.org.br/resposta-demanda}
\item Controle sintético para baseline de RD: \url{https://arxiv.org/abs/2604.18469}
\item Zhang et al.\ (2018), seq2point, AAAI-18; \url{https://github.com/nilmtk/nilmtk-contrib}
\item ANEEL, Tarifa Branca: \url{https://www.gov.br/aneel/pt-br/assuntos/tarifas/tarifa-branca}
\item Agência iNFRA, Tarifa Branca automática (2025)
\item ANEEL, tarifas homologadas: \url{https://dadosabertos.aneel.gov.br/dataset/tarifas-distribuidoras-energia-eletrica}
\item Tang et al.\ (2017), Top-K Insights, SIGMOD; Ding et al.\ (2019), QuickInsights, SIGMOD, doi:10.1145/3299869.3314037
\item PySAL esda: \url{https://pysal.org/esda/}
\item Caldas de Castro e Singer (2006), FDR em LISA, \emph{Geographical Analysis}
\item Truong et al.\ (2020), ruptures, \emph{Signal Processing}; Adams e MacKay (2007), BOCPD
\item Wickramasuriya, Athanasopoulos e Hyndman (2019), MinT, \emph{JASA} 114(526)
\item Lundberg et al.\ (2020), TreeSHAP, \emph{Nature Machine Intelligence} 2:56--67
\item Bento et al.\ (2021), TimeSHAP, KDD
\item Li (2024), GeoShapley, \emph{Annals of the AAG} 114(7)
\item Gneiting e Raftery (2007), \emph{JASA}; properscoring; xskillscore
\item Evidently: \url{https://github.com/evidentlyai/evidently}
\item alibi-detect, licença: \url{https://github.com/SeldonIO/alibi-detect/blob/master/LICENSE}
\item NannyML: \url{https://nannyml.readthedocs.io/}
\item Rabanser, Günnemann e Lipton (2019), Failing Loudly, NeurIPS
\item BIRD e Agentar-Scale-SQL: \url{https://arxiv.org/abs/2509.24403}
\item GeoBenchX: \url{https://arxiv.org/abs/2503.18129}
\item DataGemma: \url{https://arxiv.org/abs/2409.13741}
\end{enumerate}}
